\documentclass[11pt,a4paper]{coverletter}

\senderblock
  {Radheem Bin Razi}
  {Distributed Systems \& Cloud-Native Engineer}
  {Ilmenau, Germany \textbar\ \href{mailto:sheikh.radheem@gmail.com}{sheikh.radheem@gmail.com} \textbar\ \href{https://radheem.github.io/my-cv}{portfolio}}

\begin{document}

%% ===== ENGLISH =====
\recipient{Fraunhofer Institute for Digital Media Technology IDMT}{Hiring Team}{}

\opening{Dear Hiring Team,}

\vspace{8pt}\noindent\textbf{1. Why Fraunhofer IDMT?}\par\vspace{3pt}

Trustworthy AI requires more than accurate predictions; it demands systems that protect data at rest and in transit while analyzing complex signals. The Media Distribution and Security group at Fraunhofer IDMT builds exactly that kind of responsible infrastructure. I am applying for the Research/Software Engineer role to contribute to projects that balance advanced machine learning with rigorous data protection.

\vspace{8pt}\noindent\textbf{2. Why me?}\par\vspace{3pt}

My background bridges machine learning pipelines with strict security practices. I designed a containerized application that generates documents using a local language model while encrypting every output with AES-256-GCM before storage. This system also tracks every generation step through version control and quality benchmarks, ensuring reproducibility without exposing API keys. That same focus on secure, auditable workflows guided my work on an O-RAN AIML framework, where I automated the full training lifecycle using Python and Kubernetes. The project earned a top research grade and required careful handling of feature stores and model artifacts.

I also build event-driven backends that turn raw metrics into actionable insights. I configured a Kafka pipeline that routes performance data to multiple databases and visualization tools, then wrapped those flows in a Docker Compose stack for easy deployment. I am comfortable translating these backend systems into web demonstrators using Python and JavaScript, and I regularly package services with Git and Docker for consistent delivery. These experiences align directly with your need for AI software development, signal analysis support, and privacy-aware engineering.

\vspace{8pt}\noindent\textbf{3. Why now?}\par\vspace{3pt}

I am ready to transition into applied research and software engineering at a level that matches my current experience. My career has consistently moved toward building secure, data-driven systems, and Fraunhofer IDMT offers the exact environment to deepen that work. I am available to start full-time immediately, and I remain open to working-student or part-time arrangements if preferred. I am currently based in Germany and can relocate within the country within two to three weeks. A qualifying offer would only require a standard employment contract, as I am Blue Card-ready without additional sponsorship steps.

\closing{Sincerely,}{Radheem Bin Razi}

\clearpage
\selectlanguage{ngerman}

%% ===== DEUTSCH =====
\recipient{Fraunhofer Institute for Digital Media Technology IDMT}{Personalabteilung}{}

\opening{Sehr geehrtes Hiring-Team,}

\vspace{8pt}\noindent\textbf{1. Warum Fraunhofer IDMT?}\par\vspace{3pt}

Vertrauenswürdige KI erfordert mehr als präzise Vorhersagen; sie verlangt nach Systemen, die Daten im Ruhezustand und während der Übertragung schützen, während komplexe Signale analysiert werden. Die Media Distribution and Security group am Fraunhofer IDMT entwickelt genau diese Art von verantwortungsvoller Infrastruktur. Ich bewerbe mich für die Position als Research/Software Engineer, um an Projekten mitzuwirken, die fortschrittliches Machine Learning mit strengen Datenschutzstandards in Einklang bringen.

\vspace{8pt}\noindent\textbf{2. Warum ich?}\par\vspace{3pt}

Mein Hintergrund verbindet Machine-Learning-Pipelines mit strengen Sicherheitspraktiken. Ich entwickelte eine containerisierte Anwendung, die Dokumente mithilfe eines lokalen Sprachmodells generiert und jede Ausgabe vor der Speicherung mit AES-256-GCM verschlüsselt. Dieses System protokolliert zudem jeden Generierungsschritt über Versionskontrolle und Qualitätsbenchmarks, wodurch Reproduzierbarkeit gewährleistet wird, ohne API-Schlüssel preiszugeben. Dieselbe Ausrichtung auf sichere, überprüfbare Workflows leitete meine Arbeit an einem O-RAN AIML framework, bei der ich den gesamten Trainingslebenszyklus mit Python und Kubernetes automatisierte. Das Projekt erhielt eine hervorragende Bewertung und erforderte einen sorgfältigen Umgang mit Feature Stores und Model Artifacts.

Darüber hinaus entwickle ich ereignisgesteuerte Backends, die Rohmetriken in umsetzbare Erkenntnisse umwandeln. Ich konfigurierte eine Kafka pipeline, die Leistungsdaten an mehrere Datenbanken und Visualisierungstools weiterleitet, und verpackte diese Datenflüsse anschließend in einen Docker Compose stack für eine einfache Bereitstellung. Ich bin erfahren darin, diese Backend-Systeme mit Python und JavaScript in Web-Demonstratoren zu übersetzen, und verpacke Dienste regelmäßig mit Git und Docker für eine konsistente Auslieferung. Diese Erfahrungen entsprechen direkt Ihren Anforderungen an KI-Softwareentwicklung, Unterstützung bei der Signalanalyse und datenschutzorientierte Ingenieursarbeit.

\vspace{8pt}\noindent\textbf{3. Warum jetzt?}\par\vspace{3pt}

Ich bin bereit, in die angewandte Forschung und Softwareentwicklung in einem Umfang einzusteigen, der meiner aktuellen Erfahrung entspricht. Mein Karriereweg war stets darauf ausgerichtet, sichere, datengetriebene Systeme aufzubauen, und das Fraunhofer IDMT bietet genau das Umfeld, um diese Arbeit zu vertiefen. Ich stehe ab sofort in Vollzeit zur Verfügung und bin auch an Werkstudentenmodellen oder Teilzeitregelungen interessiert, falls dies bevorzugt wird. Ich befinde mich derzeit in Deutschland und kann innerhalb von zwei bis drei Wochen innerhalb des Landes umziehen. Ein entsprechendes Angebot würde lediglich einen regulären Arbeitsvertrag erfordern, da ich bereits Blue Card-ready bin und keine zusätzlichen Sponsorship steps notwendig sind.

\closing{Mit freundlichen Grüßen,}{Radheem Bin Razi}

\end{document}
